\documentclass{article}

\usepackage{booktabs}
\usepackage{multirow}
\usepackage{amsmath}
\usepackage{hyperref}
\usepackage{overpic}
\usepackage{amssymb}

\usepackage[accepted]{dlaiml2025}

%%% STUDENTS: FILL IN WITH YOUR OWN INFORMATION
\dlaititlerunning{Knowledge Distillation from Ensemble Teachers to a Tiny MLP Student}

\begin{document}

\twocolumn[
%%% STUDENTS: FILL IN WITH YOUR OWN INFORMATION
\dlaititle{Knowledge Distillation from Ensemble Teachers to a Tiny MLP Student}

\begin{center}\today\end{center}

\begin{dlaiauthorlist}
%%% STUDENTS: FILL IN WITH YOUR OWN INFORMATION
\dlaiauthor{Eliseo Valerio Gagliardi}{}
\end{dlaiauthorlist}

%%% STUDENTS: FILL IN WITH YOUR OWN INFORMATION
\dlaicorrespondingauthor{Eliseo Valerio Gagliardi}{gagliardi.2070803@studenti.uniroma1.it}

\vskip 0.3in
]

\printAffiliationsAndNotice{}

\begin{abstract}
%
This project investigates the transfer of information from an ensemble of machine learning teachers (Random Forest, SVM, KNN) to a compact MLP student. The goal is to maintain accuracy on MNIST while reducing computational demands, offering a pathway for efficient deployment of neural networks on
constrained hardware. 
\cite{hinton2015distilling}
\paragraph*{Code.} 
You can find the project as a notebook in Google Colab in this Github repository: \url{https://github.com/Spiky03/ML-project-2025}.
%
\end{abstract}

% ------------------------------------------------------------------------------
\section{Background and Motivation}

Modern AI systems often require substantial computing power and energy, creating obstacles to real-world deployment on mobile and embedded devices. Knowledge distillation addresses this by enabling smaller models to learn from larger, more accurate ones. Unlike conventional studies that rely on deep neural networks as teachers, this project demonstrates that ensembles of traditional machine learning models can also serve as powerful knowledge sources. This exploration widens the scope of distillation, showing how cross-paradigm approaches can yield compact yet effective models.



% ------------------------------------------------------------------------------
\section{Implementation}

The team used the MNIST dataset, reduced to 10,000 samples for computational efficiency. The teacher ensemble consisted of three diverse models: Random Forest, SVM, and KNN. These models generated calibrated probability distributions, which were distilled into a three-layer Tiny MLP. Training combined standard cross-entropy loss with softened teacher outputs, allowing the student to absorb richer supervision signals beyond hard labels. \cite{romero2015fitnets}

\paragraph{Distillation loss.} The notebook defines the distillation loss as 
\begin{equation} 
    \mathcal{L} = \alpha,\mathcal{H}(y, p_s) + (1-\alpha),T^{2},\mathrm{KL}!\big(p_t^{(T)} ,|, p_s^{(T)}\big), 
\end{equation} 
 aligns the student distribution $p_s^{(T)}$ with the softened teacher output $p_t^{(T)}$ at temperature $T$. Teacher probabilities are stabilized using \texttt{softmax(log(teacher\_probs)/T)}, and the KL divergence is computed via \texttt{F.kl\_div}. Note that, unlike the conventional formulation, the notebook assigns $\alpha$ to the cross-entropy term and $(1-\alpha)$ to the distillation term, which is simply a choice of convention.

 

% ------------------------------------------------------------------------------
\section{Results and Insights}\label{sec:results}

\begin{figure}[b]
    \centering
    \begin{overpic}[width=0.99\linewidth]{./Loss_acc.png}
    \end{overpic}
\end{figure}

The distilled Tiny MLP outperformed its baseline counterpart trained solely on labels, achieving better validation accuracy. This confirms that ensembles of relatively simple models can act as effective collective teachers. Importantly, the approach showcases a scalable strategy for deploying accurate AI on constrained hardware without relying on massive deep learning architectures.

\begin{table}[h!]
\caption{Performance comparison.}
\label{tab:results}
\begin{center}
\begin{small}
\begin{tabular}{p{0.16\linewidth} | ccccc}
\toprule
& \multirow{2}{0.1\linewidth}{Random Forest}& \multirow{2}{0.2\linewidth}{SVM (Calibrated)}& \multirow{2}{0.1\linewidth}{KNN}& \multirow{2}{0.2\linewidth}{Ensemble}& \multirow{2}{0.1\linewidth}{Student (KD)}\\
Models\\
\midrule
Accuracy      & 94.87\% & 95.82.0\% &  94.34\% & 96.09\% &  95.30\%  \\
\bottomrule
\end{tabular}
\end{small}
\end{center}
\end{table}

An example that resonates with this concept is \textbf{DeepSeek}, which revolutionized the AI market by proving that efficiency-driven design choices can outperform raw scale. Just as DeepSeek optimized inference speed and cost while maintaining state-of-the-art performance, this project illustrates that knowledge distillation can unlock similar efficiencies for compact models.

\bibliography{references.bib}
\bibliographystyle{dlaiml2025}

\end{document}

